\subsection{A common head at a general center}

We now work in the finite-height forest space $X=J^*$. For each head
$H$, write $P=P_H$ and $Q=Q_H$. The next two
applications choose the head before the displacement. Their center and
radius assumptions differ, but they use the same original-coefficient
mass estimate. For a finite atom law tested by $p=P_Hp$ with $\rho(p)\le1$, put
\[
 k=\mathbb E L(A),\qquad
 \kappa=\|((\mathbb EA)(e_g))_{g\in G_H}\|_2,\qquad
 M_c=(\mathbb E|c_g(A)|)_{g\in G_H}.
\]
The gate-coordinate identity gives
$\beta=\|\mathbb Ec\|_2\le\kappa+\sqrt{2k}$. Hence the original-product
bound in Lemma~\ref{stab:coef:gate-bound} gives
\[
 \|M_c\|_2\le
 \sqrt{(\kappa+\sqrt{2k})^2+4k+4\sqrt{2k}}.
\]
For an equal mixture, each conditional crossing tail costs at most
$2\|M_c\|_2$, and its wholly exterior part costs at most
$2\sqrt{2k}$. These follow from
Lemma~\ref{stab:coef:tail-assembly} and the nonnegative pooled
exterior budget. The proofs below specify $k$, $\kappa$, and the
endpoint-approximation error in their respective settings.

\subsubsection{Unit centers and an enlarged lens}
\label{molb:sec:original-coefficients}

\begin{proposition}\label{molb:prop:original-lens}
For \(0<\delta<1\) and every unit \(x\in X\), one can choose a finite
ancestor-closed \(H\) such that \(\|x\pm y\|\le1+\delta\) implies
\[
 \|Qy\|\le\eta(\delta):=(1+\delta)
 \left\{2\sqrt{(2\delta+\sqrt{6\delta})^2
                    +4(3\delta+\sqrt{6\delta})}
             +2\sqrt{6\delta}+\delta\right\}.
\]
In particular \(\eta(\delta)\to0\) as \(\delta\downarrow0\).
\end{proposition}

\begin{proof}
Choose \(H\) so that \(\|Q_Hx\|\le\delta\), and choose a primal
unit test \(p=P_Hp\) with \(x(p)\ge1-\delta\). Such a test norms
\(P_Hx\), whose norm is at least \(1-\delta>0\).
Both choices depend only on \(x,\delta\), before \(y\) is given.

Approximate each scaled endpoint \((x\pm y)/(1+\delta)\) within
\(\delta\) by a finite atom mean, and mix the two laws equally.
With \(k=\mathbb E(1-A(p))\), the pooled evaluation and tail satisfy
\[
 k\le1-\frac{1-\delta}{1+\delta}+\delta\le3\delta,
 \qquad
 \|Q_H\mathbb EA\|\le\frac{\|Q_Hx\|}{1+\delta}+\delta
 \le2\delta.
\]
Thus the pooled gate norm \(\kappa\) is at most \(2\delta\).
The product-transfer estimate used here is
\begin{equation}\label{molb:eq:original-products}
 \sum_g(-c_g(A)c_g(A'))_+
 \le L(A)+L(A')+\sqrt{2L(A)}+\sqrt{2L(A')},
\end{equation}
proved in Lemma~\ref{stab:coef:original-products} by the
one-Lipschitz property of the positive part and the head-conflict bound.
Applying the displayed mass estimate above gives
\[
 \|M_c\|_2\le
 B:=\sqrt{(2\delta+\sqrt{2k})^2+4(k+\sqrt{2k})}.
\]
For either conditional law, the crossing tail has norm at most \(2B\)
and the wholly exterior part has norm at most \(2\sqrt{2k}\).
Restore the endpoint approximation error \(\delta\), take half the
difference of the two endpoint tails, and multiply by \(1+\delta\).
The bound is therefore
\[
 \|Q_Hy\|\le(1+\delta)(2B+2\sqrt{2k}+\delta).
\]
Using \(k\le3\delta\) gives exactly the stated \(\eta(\delta)\).
Its leading order is \(\delta^{1/4}\), inherited from the
original-product estimate.
\end{proof}

\subsubsection{Near-unit centers in the unit ball}
\label{molx:actual-center-section}

Here the endpoints lie in the unit ball and the center is only
assumed nearly unit. The finite head must again be fixed uniformly for
all displacements in its lens.

\begin{proposition}\label{molx:actual-center-fourth}
For $0<\lambda<1$ and $x\in B_X$ with $\|x\|\ge1-\lambda$,
there is a finite ancestral set $H$, depending only on $x$ and
$\lambda$, such that every $y\in X$ with $x\pm y\in B_X$ satisfies
\begin{equation}\label{molx:actual-center-fourth-bound}
 \|Q_Hy\|\le G_{\mathrm{orig}}(\lambda)
 :=\lambda+2\sqrt{6\lambda}
 +2\sqrt{(2\lambda+\sqrt{6\lambda})^2
                  +12\lambda+4\sqrt{6\lambda}}.
\end{equation}
In particular $G_{\mathrm{orig}}(\lambda)\to0$ as $\lambda\downarrow0$.
\end{proposition}

\begin{proof}
By density of finitely supported vectors in $J$, choose
$p\in B_J$ with finite support and $x(p)\ge\|x\|-\lambda$.
Use the contractive exhaustion in Lemma~\ref{lem:forest-structure}
to choose a finite ancestral $H$ containing this support and satisfying
$\|Q_Hx\|\le\lambda$.  Make these choices before selecting $y$.

By Corollary~\ref{cor:closed-atomic-ball}, choose finite represented
atom laws with means $z_+,z_-$ such that
$\|z_\pm-(x\pm y)\|\le\lambda$.  Write $\mathbb E_\pm$ for
their expectations and $\mathbb E=(\mathbb E_++\mathbb E_-)/2$.
For $z=(z_++z_-)/2$ one has
\begin{equation}\label{molx:actual-center-budget}
 \|Q_Hz\|\le2\lambda,\qquad
 z(p)\ge1-3\lambda.
\end{equation}

Use the notation of Section~\ref{stab:toolkit}.  For an atom
$A=\sum_i c_i d_{S_i}$, put $s_i=p(S_i)$ and $L=1-A(p)$.
Its wholly exterior part is $O(A)$, and $c_g,s_g$ are its original
and corrected crossing coefficients at $g\in G_H$.  The pointwise
budgets in Lemma~\ref{stab:coef:pointwise} give
\[
 \|O(A)\|\le\sqrt{2L},\qquad
 \|c-s\|_2\le\sqrt{2L},\qquad
 \|c\|_2,\|s\|_2\le1.
\]
Here the last two vectors are the crossing arrays.  In particular
$\mathbb EL\le3\lambda$ and
$\mathbb E\sqrt{2L}\le\sqrt{6\lambda}$.

For two independent represented atoms, the original-product estimate
of Lemma~\ref{stab:coef:original-products} reads
\begin{equation}\label{molx:actual-center-transfer}
 \sum_g(-c_gc'_g)_+
 \le L+L'+\sqrt{2L}+\sqrt{2L'}.
\end{equation}
Its head comparison may use either the truncation or the two-piece
proof in Lemma~\ref{stab:coef:head-conflict}; both give the same
transfer bound. This is the estimate underlying the common mass
bound above.
Thus, for $a_g=\mathbb E(c_g)_+$ and
$b_g=\mathbb E(-c_g)_+$, independence yields
\[
 4\sum_g a_gb_g\le12\lambda+4\sqrt{6\lambda}.
\]
The gate-coordinate identity
\eqref{stab:coef:gate-coordinates} also gives
\[
 \|a-b\|_2=\|\mathbb Ec\|_2
 \le\|Q_Hz\|+\mathbb E\|O(A)\|
 \le2\lambda+\sqrt{6\lambda}.
\]
Consequently the original absolute mass $m=a+b=\mathbb E|c|$
satisfies
\begin{equation}\label{molx:actual-center-mass}
 \|m\|_2^2\le(2\lambda+\sqrt{6\lambda})^2
                         +12\lambda+4\sqrt{6\lambda}.
\end{equation}

For either endpoint law, positivity of the mixture weights gives
\[
 \mathbb E_\pm\|O(A)\|\le2\sqrt{6\lambda},\qquad
 \mathbb E_\pm|c_g|\le2m_g.
\]
Lemma~\ref{stab:coef:tail-assembly}, or its proof by maximizing
one test segment in each cone, therefore bounds the whole crossing
tail by $2\|m\|_2$.  The segments selected in different cones
are disjoint, so their primal test sums have squared sum at most
one; this is what turns the absolute masses into the whole-tail
bound.  Thus
$\|Q_Hz_\pm\|\le2\sqrt{6\lambda}+2\|m\|_2$.
Finally
$\|y-(z_+-z_-)/2\|\le\lambda$, and hence
\[
 \|Q_Hy\|\le\lambda+
 \tfrac12\bigl(\|Q_Hz_+\|+\|Q_Hz_-\|\bigr).
\]
Substitute \eqref{molx:actual-center-mass}.  All choices of $H$
and $p$ preceded $y$, as required.
\end{proof}


\subsection{Midpoint consequences with their original losses}
\label{stab:coef:moduli}

We return to either coordinate model $Z$ and use the averaged modulus defined in \eqref{eq:averaged-modulus},
with perturbations of norm at least one. The following conversions
retain their different lower bounds on the endpoint norms.

\begin{proposition}\label{stab:coef:modulus-conversions}
For $0<t<1$, each of the following conclusions holds.
\begin{enumerate}
\item If $e>0$ satisfies
\[
 (\sqrt{12}+2\sqrt2)\sqrt{(1+e)e}<t/2,
\]
then $\widehat\delta_Z(t)\ge e/4$.
\item If $0<c<1$, $R=1+4c$, and $R\omega_{\mathrm{coef}}(5c/R)<t$, then
$\widehat\delta_Z(t)\ge c$.
\item If $d>0$, $3d<1$, and $(1+2d)\omega_{\mathrm{coef}}(3d)<t$, then
$\widehat\delta_Z(t)\ge d/2$.
\item The completed atomic-ball route gives
$\widehat\delta_Z(t)\ge t^2/16384$.
\end{enumerate}
\end{proposition}

\begin{proof}
For the first assertion, approximate any unit $x$ by a unit finite
head $x_0=P_Hx_0$ with $\|x-x_0\|<e/4$, and use
$F=\ker P_H$. For $y\in F$ with $\|y\|\ge1$,
Corollary~\ref{stab:coef:unit-head-lens} rules out both
$\|x_0\pm ty\|\le1+e$, since their lens bound is less than
$t/2<t\|y\|$. Both endpoint norms are at least one by projection.
Their average is therefore greater than $1+e/2$, and replacing
$x_0$ by $x$ leaves it greater than $1+e/4$.

For the second assertion choose $H$ with
$\|x-P_Hx\|\le c$ and again let $F=\ker P_H$. Both
$\|x\pm ty\|\ge\|P_Hx\|\ge1-c$. If their average were
at most $1+c$, each would be at most $1+3c$. It would follow
that $\|P_Hx\pm ty\|\le1+4c=R$. For every finite ancestral
$K\supset H$, apply \eqref{stab:coef:fourth-root} to
$P_Hx/R$ and $tP_Ky/R$. Its deficit is at most
$1-(1-c)/R=5c/R$, so
\[
 t\|P_Ky\|\le R\omega_{\mathrm{coef}}(5c/R)<t.
\]
Exhaustion gives $t\|y\|<t$, a contradiction. Notice that the
definition of $\omega_{\mathrm{coef}}$ for all nonnegative arguments
makes this parameter substitution meaningful even before $c$ is made small.

For the third assertion choose $H$ with
$\|x-P_Hx\|<d$, and let $f\in Z^*$ have norm one and
$f(x)=1$. Put $F=\ker P_H\cap\ker f$. Then
$\|x\pm ty\|\ge1$ for $y\in F$. If both were at most
$1+d$, then $\|P_Hx\pm ty\|\le1+2d$, and the normalized
head would have norm at least
\[
 \frac{1-d}{1+2d}\ge1-3d.
\]
Proposition~\ref{stab:coef:exact-tail} would give
$t\le(1+2d)\omega_{\mathrm{coef}}(3d)$ for $\|y\|\ge1$, contrary to the
chosen parameter. Thus the average is greater than $1+d/2$.

Finally set $d=(t/64)^2$ for the fourth assertion. Choose $H$
so that $\|x-P_Hx\|<d/8$. The normalized head
$x_0=P_Hx/\|P_Hx\|$ then exists and satisfies
$\|x-x_0\|<d/4$, because normalization costs at most the
projection error again. For $y\in\ker P_H$ of norm at least
one, Proposition~\ref{stab:coef:approximate-tail} and
$32\sqrt d=t/2<t$ show that one endpoint $x_0\pm ty$ has
norm greater than $1+d$. Both have norm at least one. Their
average exceeds $1+d/2$, and replacing $x_0$ by $x$ leaves
more than $1+d/4=1+t^2/16384$.
\end{proof}

\subsection{A conditional estimate using prefix truncation}
\label{stab:coef:conditional}

One can also keep an endpoint law separate while using the head-sign
comparison. The following argument retains the wholly exterior term
under that law; its coefficient is smaller than the one obtained by
bounding every conditional term through the pooled expectation.

\begin{proposition}\label{stab:coef:conditional-modulus}
For $0<t<1$, set $\delta=(t/25)^2$. If $x_0=P_Hx_0$ and $y$
are unit vectors, $P_Hy=0$, and both are finitely supported, then
\begin{equation}\label{stab:coef:finite-strict}
 \max\{\|x_0+ty\|,\|x_0-ty\|\}>1+\delta.
\end{equation}
For general unit $y\in\ker P_H$, the same inequality holds with
$\ge$ in place of $>$. Moreover,
\begin{equation}\label{stab:coef:conditional-modulus-eq}
 \widehat\delta_Z(t)\ge t^2/2500.
\end{equation}
\end{proposition}

\begin{proof}
If \eqref{stab:coef:finite-strict} failed, represent the two scaled
endpoints by exact finite atom laws and use their equal mixture as in
\eqref{stab:coef:exact-laws}. Put
$\eta=\delta/(1+\delta)$; this is the deficit of the scaled
unit head. The conditional and pooled deficits all equal $\eta$.
Using the \emph{truncation} proof of
Lemma~\ref{stab:coef:head-conflict}, the calculations
\eqref{stab:coef:exact-means}--\eqref{stab:coef:twelve} give
\[
 \|\mathbb Ec\|_2\le\sqrt{2\eta},\quad
 \|\mathbb Es\|_2\le2\sqrt{2\eta},\quad
 \|M_s\|_2^2\le12\eta,\quad
 \|M_c\|_2\le\sqrt{12\eta}+\sqrt{2\eta}.
\]
In the plus law, $\mathbb E_+\|O(A)\|\le\sqrt{2\eta}$
because its own mean deficit is $\eta$. For the crossing part use
the disjoint-test-segment proof of
Lemma~\ref{stab:coef:tail-assembly} and the bound of its conditional
absolute mean by twice the pooled one. Therefore
\begin{equation}\label{stab:coef:twelve-conditional}
 \frac{t}{1+\delta}
 \le2(\sqrt{12\eta}+\sqrt{2\eta})+\sqrt{2\eta}
 \le12\sqrt\delta.
\end{equation}
Indeed $2\sqrt{12}+3\sqrt2<12$ and $\eta\le\delta$.
But $\delta<1/625$ and
\[
 12(1+\delta)\sqrt\delta
 =\frac{12}{25}(1+\delta)t<t,
\]
contradicting \eqref{stab:coef:twelve-conditional}.

For a completed unit tail $y$, take finite ancestral projections
$P_Ky\to y$ with $K\supset H$. For all sufficiently large $K$
the projections are nonzero. Normalize them to unit vectors;
they still vanish on $H$ and converge to $y$. Passing to the limit
in the finite strict inequality yields the stated non-strict one.

Given an arbitrary unit center $x$, choose $H$ so that its normalized
head $x_0$ satisfies $\|x-x_0\|<\delta/2$. Choose a norm-one
$f\in Z^*$ with $f(x)=1$, and set
$F=\ker P_H\cap\ker f$. For unit $y\in F$, both actual
endpoint norms are at least one, and the maximum is at least
$1+\delta/2$. Their average is at least $1+\delta/4$.
For a fixed direction, the average
$s\mapsto(\|x+sy\|+\|x-sy\|)/2$ is even and convex, hence
nondecreasing for $s\ge0$. This extends the bound to all
$y\in F$ of norm at least one and gives
$\widehat\delta_Z(t)\ge\delta/4=t^2/2500$.
\end{proof}

\subsection{Separated families in a lens}

For the forest space $X$, the function $g(s)$ of Proposition~\ref{mola:actual} tends to zero and is of order $s^{1/4}$ as
$s\downarrow0$.  To see the separated-lens consequence directly,
fix $t>0$ and choose $0<\Delta<1/2$ such that
\[
 (1+2\Delta)g(4\Delta)+\Delta<t.
\]
If an infinite family $(y_j)$ satisfies $\|x\|>1-\Delta$,
$\|x\pm y_j\|\le1$, and $\|y_i-y_j\|\ge2t$ for distinct indices,
choose a finite head with
$\|Q_Hx\|\le\Delta$ and a pair with
$\|P_H(y_i-y_j)/2\|\le\Delta$.  Put $v=(y_i-y_j)/2$.
The vectors $x\pm v$ are averages of opposite endpoints and lie in $B_X$.
The vectors $P_Hx/(1+2\Delta)$ and $Q_Hv/(1+2\Delta)$ satisfy the
hypotheses of Proposition~\ref{mola:actual}; the first has deficit
at most $4\Delta$, whereas the second has norm at least
$(t-\Delta)/(1+2\Delta)$.  This contradicts the choice of $\Delta$.

The actual-center estimate gives a direct version with one head for the whole family.
For example, choose $\lambda\in(0,1)$ so that
$2G_{\mathrm{orig}}(\lambda)<\varepsilon/2$.  If $x\pm y_j\in B_X$ for an
infinite $\varepsilon$-separated family and
$\|x\|\ge1-\lambda$, the same head works for every $j$.
The bounded sequence $(P_Hy_j)$ has two terms at distance less
than $\varepsilon/2$, while the corresponding tail difference has
norm at most $2G_{\mathrm{orig}}(\lambda)<\varepsilon/2$, a contradiction.
